% ============================================================
% OP_GUT3b_D1_blocks_v2.tex  (EXACT INSERTION BLOCKS for review;
% NOT yet implemented). v2 = mandatory branch-(iii) lattice
% correction per GPT review of v3/D1-v1: combined residual
% lattice is c in 2Z with lambda + c = 1 (mod 6); c in 6Z is the
% pure B-L shift at lambda = 1 only. Scope per GPT verdicts:
% compact block "Hypercharge normalisation: distinct notions and
% lattice interlock" + one-line OP-GUT3a/3b registry echoes.
% Corrected sign Q = -t/3 (mod 1) throughout; mandatory framing
% sentence included; recorded global-structure citation reused
% (narrow consistency claim only - no Gamma selection, no
% line-operator claims); N3 (g', Z_B, sin^2 theta_W) absent.
% ============================================================

% ---------- ANCHOR 1 ----------
% Insert AFTER (end of the anomaly solution-structure block):
%   "Its whole-framework consistency status is recorded in the global
%    Level-4 self-consistency checklist~(\S\ref{sec:global_level4_checklist})."
% and BEFORE:
%   "\subsubsection*{P7-restricted matter audit (OP-GUT2): bounded
%    first pass}"

\paragraph{Hypercharge normalisation: distinct notions and
lattice interlock.}
Four notions must be kept separate: the rational pattern of
hypercharge ratios within a generation (the core of OP-GUT3b);
the charge-lattice/centre structure (OP-GUT3a and its extension
here); the coupling magnitude ($g'$, $Z_B^{(0)}$,
$\sin^2\theta_W$ --- OP-GUT4/OP-GUT7, not part of OP-GUT3b); and
the bookkeeping convention $Q=T_3+Y/2$ used throughout.
For a multiplet with triality $t$ ($\mathbf3\to+1$,
$\bar{\mathbf3}\to-1$, colour singlet $\to0$), duality $d$ (weak
doublet $\to1$, singlet $\to0$) and hypercharge $Y$, the centre
element $(\omega\mathbf 1_3,\,-\mathbf 1_2,\,e^{i\pi Y})$,
$\omega=e^{2\pi i/3}$, acts trivially iff the interlock
congruence
\begin{equation}
  \frac{2t}{3}+d+Y\;\in\;2\mathbb Z
  \quad\Longleftrightarrow\quad
  3Y+3d+2t\equiv0\ (\mathrm{mod}\ 6)
  \label{eq:interlock_congruence}
\end{equation}
holds.
All recorded one-generation assignments~\eqref{eq:hypercharges}
satisfy it (values $2,-2,0,0,2,0$), so the inherited content is
consistent with the maximal quotient
$(SU(3)\times SU(2)\times U(1)_Y)/\mathbb Z_6$; consistency does
not select $\Gamma$, and distinguishing the
$\mathbb Z_6/\mathbb Z_3/\mathbb Z_2$/trivial options requires
extended-object data~\cite{AharonyEtAl2013}.
For colour singlets in the recorded singlet/doublet sectors
($d=2T\ \mathrm{mod}\ 2$), the congruence reduces to the recorded
electroweak $\mathbb Z_2$ condition $2T+Y\in2\mathbb Z$
\eqref{eq:charge_lattice}; it is the colour-completed extension
of that refinement.
With $Q=T_3+Y/2$ it implies $Q\equiv-t/3\ (\mathrm{mod}\ 1)$:
colour singlets carry integer electric charge, triplets
$Q\equiv-\tfrac13$ and antitriplets
$Q\equiv+\tfrac13\ (\mathrm{mod}\ 1)$ --- the rule correlating
allowed electric classes with colour representations.
If the observed quark charges are admitted as input, the minimal
electric lattice compatible with them has unit $e/3$; the
nontrivial content is the correlation rule, not the arithmetic
unit, and for the colour completion (OP-GUT1) the rule is a
requirement: the emergent colour module must realise the
$\mathbb Z_3$ centre with exactly this triality--charge
correlation.
The congruence also acts on the residual freedoms of the
recorded anomaly solution structure.
In the branches without a free-hypercharge $\nu^c$, the
rescaling $Y\to\lambda Y_{\rm SM}$ is discretised: the quark
doublet gives $\lambda\equiv1\ (\mathrm{mod}\ 6)$, and the
remaining multiplets impose weaker congruences
($\lambda\equiv1\ \mathrm{mod}\ 3$ or
$\lambda\equiv1\ \mathrm{mod}\ 2$), so the combined condition is
$\lambda\equiv1\ (\mathrm{mod}\ 6)$.
In the free-hypercharge $\nu^c$ branch, for the pure $B{-}L$
shift around the observed representative ($\lambda=1$) the
interlock requires $c\in6\mathbb Z$.
For the combined deformation
$Y'=\lambda Y_{\rm SM}+c\,(B{-}L)$, however, the residual
lattice is
\[
  c\in2\mathbb Z,\qquad \lambda+c\equiv1\pmod 6,
\]
equivalently $c=2m$, $\lambda=1-2m+6n$ with $m,n\in\mathbb Z$
(the point $(\lambda,c)=(-1,2)$ realises the recorded
$u^c\!\leftrightarrow\!d^c$ relabelling).
Thus the interlock discretises the $B{-}L$ flat direction but
does not eliminate it; the observed assignment remains the
representative $(\lambda,c)=(1,0)$, and selecting it is still an
additional minimality/completion choice
(cf.\ the minimality convention of the matter audit below).
This is a consistency requirement on completions, not a
derivation of the observed hypercharges.
OP-GUT3a is extended at the congruence level; OP-GUT3b remains
Open, sharpened to selection within the explicitly characterised
discrete residual family.

% ---------- ANCHOR 2a (registry echo, tab:op_gut) ----------
% In the OP-GUT3a row, REPLACE:
%   "    EW $\mathbb{Z}_2$ refinement conditional~(B)"
% WITH:
%   "    EW $\mathbb{Z}_2$ refinement conditional~(B);
%        colour-interlock congruence recorded
%        (\S\ref{sec:anomaly_architecture})"

% ---------- ANCHOR 2b (registry echo, tab:op_gut) ----------
% In the OP-GUT3b row, REPLACE:
%   "  & Derive observed charge assignments of quarks and leptons
%      & Open \\"
% WITH:
%   "  & Derive observed charge assignments of quarks and leptons;
%       sharpened to selection within an explicitly characterised
%       discrete residual family (\S\ref{sec:anomaly_architecture})
%      & Open \\"
